% Part: first-order-logic
% Chapter: sequent-calculus
% Section: proving-things

\documentclass[../../../include/open-logic-section]{subfiles}

\begin{document}

\iftag{FOL}
      {\olfileid[pt]{fol}{seq}{pro}}
      {\olfileid[pt]{pl}{seq}{pro}}

\olsection{Exemplos de \usetoken{P}{derivation}}

\begin{ex}
Dê uma $\Log{LK}$-!!{derivation} para o sequente $!A \land !B \Sequent !A$.

Começamos escrevendo o sequente final desejado na base da
!!{derivation}.
\begin{prooftree}
\AxiomC{}
\UnaryInf$!A\land !B \fCenter !A$
\end{prooftree}
Em seguida, precisamos descobrir que tipo de inferência poderia ter um
sequente inferior dessa forma. Poderia ser uma regra estrutural, mas é uma
boa ideia começar procurando uma regra lógica. O único conectivo
lógico que ocorre no sequente inferior é $\land$,
então procuramos uma regra de $\land$ e, como o símbolo $\land$
ocorre no antecedente, estamos diante da regra \LeftR{\land}.
\begin{prooftree}
\AxiomC{}
\RightLabel{\LeftR{\land}}
\UnaryInf$!A\land !B \fCenter !A$
\end{prooftree}
Há duas opções para o que poderia ter sido o sequente superior da
inferência \LeftR{\land}: poderíamos ter um sequente superior $!A
\Sequent !A$ ou $!B \Sequent !A$. Claramente, $!A \Sequent !A$ é um
sequente inicial (o que é bom), enquanto $!B \Sequent !A$ não é
derivável em geral. Preenchemos o sequente superior:
\begin{prooftree}
\Axiom$!A \fCenter !A$
\RightLabel{\LeftR{\land}}
\UnaryInf$!A\land !B \fCenter !A$
\end{prooftree}
Agora temos uma $\Log{LK}$-!!{derivation} correta do sequente $!A
\land !B \Sequent !A$.
\end{ex}

\begin{ex}
Dê uma $\Log{LK}$-!!{derivation} para o sequente $\lnot !A \lor !B
\Sequent !A \lif !B$.

Comece escrevendo o sequente final desejado na base da !!{derivation}.
\begin{prooftree}
\AxiomC{}
\UnaryInf$\lnot !A \lor !B \fCenter !A \lif !B$
\end{prooftree}
Para encontrar uma regra lógica que pudesse nos dar esse sequente final,
olhamos para os conectivos lógicos no sequente final: $\lnot$, $\lor$ e
$\lif$. No momento, só nos importam $\lor$ e $\lif$, porque
são !!{main operator}s de !!{sentence}s no sequente final,
enquanto $\lnot$ está dentro do escopo de outro conectivo, de modo que
cuidaremos dele depois. Nossas opções de regras lógicas para a inferência
final são, portanto, a regra \LeftR{\lor} e a regra \RightR{\lif}.
Poderíamos escolher qualquer uma das regras, na verdade, mas vamos escolher
a regra \RightR{\lif} (quando não fosse por outro motivo, porque ela nos
permite adiar a divisão em dois ramos). De acordo com a forma das
inferências \RightR{\lif} que podem produzir o sequente inferior, isto deve
ser assim:
\begin{prooftree}
\AxiomC{}
\UnaryInf$ !A, \lnot !A \lor !B \fCenter !B $
\RightLabel{\RightR{\lif}} \UnaryInf$ \lnot !A \lor !B \fCenter !A \lif !B $
\end{prooftree}

Se movermos $\lnot !A \lor !B$ para a parte externa do antecedente, podemos
aplicar a regra \LeftR{\lor}. De acordo com o esquema, isto deve se dividir
em dois sequentes superiores como segue:
\begin{prooftree}
\AxiomC{}
\UnaryInf$\lnot !A, !A \fCenter !B$
\AxiomC{}
\UnaryInf$!B, !A \fCenter !B$
\RightLabel{\LeftR{\lor}}
\BinaryInf$ \lnot !A \lor !B, !A \fCenter !B $
\RightLabel{\RightR{\Exchange}}
\UnaryInf$ !A, \lnot !A \lor !B \fCenter !B $
\RightLabel{\RightR{\lif}}
\UnaryInf$ \lnot !A \lor !B \fCenter !A \lif !B $
\end{prooftree}
Lembre-se de que estamos tentando subir até sequentes iniciais; parece
que estamos bem perto!{} O ramo direito está a apenas um enfraquecimento e
uma troca de um sequente inicial, e então estará pronto:
\begin{prooftree}
\AxiomC{}
\UnaryInf$\lnot !A, !A \fCenter !B$
\Axiom$!B \fCenter !B$
\RightLabel{\LeftR{\Weakening}}
\UnaryInf$!A, !B \fCenter !B$
\RightLabel{\LeftR{\Exchange}}
\UnaryInf$!B, !A \fCenter !B$
\RightLabel{\LeftR{\lor}}
\BinaryInf$\lnot !A \lor !B, !A \fCenter !B $
\RightLabel{\RightR{\Exchange}}
\UnaryInf$ !A, \lnot !A \lor !B \fCenter !B $
\RightLabel{\RightR{\lif}}
\UnaryInf$ \lnot !A \lor !B \fCenter !A \lif !B $
\end{prooftree}

Agora, olhando para o ramo esquerdo, o único conectivo lógico em qualquer
!!{sentence} é o símbolo $\lnot$ nas !!{sentence}s do antecedente, então
estamos diante de uma instância da regra \LeftR{\lnot}.
\begin{prooftree}
\AxiomC{}
\UnaryInf$ !A \fCenter !B, !A$
\RightLabel{\LeftR{\lnot}}
\UnaryInf$\lnot !A, !A \fCenter !B$
\Axiom$!B \fCenter !B$
\RightLabel{\LeftR{\Weakening}}
\UnaryInf$!A, !B \fCenter !B$
\RightLabel{\LeftR{\Exchange}}
\UnaryInf$!B, !A \fCenter !B$
\RightLabel{\LeftR{\lor}}
\BinaryInf$\lnot !A \lor !B, !A \fCenter !B $
\RightLabel{\RightR{\Exchange}}
\UnaryInf$ !A, \lnot !A \lor !B \fCenter !B $
\RightLabel{\RightR{\lif}}
\UnaryInf$ \lnot !A \lor !B \fCenter !A \lif !B $
\end{prooftree}
De modo semelhante a como finalizamos o ramo direito, estamos a apenas um
enfraquecimento e uma troca de finalizar também este ramo esquerdo.
\begin{prooftree}
\Axiom$!A \fCenter !A$
\RightLabel{\RightR{\Weakening}}
\UnaryInf$ !A \fCenter !A, !B$
\RightLabel{\RightR{\Exchange}}
\UnaryInf$ !A \fCenter !B, !A$
\RightLabel{\LeftR{\lnot}}
\UnaryInf$\lnot !A, !A \fCenter !B$
\Axiom$!B \fCenter !B$
\RightLabel{\LeftR{\Weakening}}
\UnaryInf$!A, !B \fCenter !B$
\RightLabel{\LeftR{\Exchange}}
\UnaryInf$!B, !A \fCenter !B$
\RightLabel{\LeftR{\lor}}
\BinaryInf$\lnot !A \lor !B, !A \fCenter !B $
\RightLabel{\RightR{\Exchange}}
\UnaryInf$ !A, \lnot !A \lor !B \fCenter !B $
\RightLabel{\RightR{\lif}}
\UnaryInf$ \lnot !A \lor !B \fCenter !A \lif !B $
\end{prooftree}
\end{ex}

\begin{ex}
Dê uma $\Log{LK}$-!!{derivation} do sequente $\lnot !A \lor \lnot !B
\Sequent \lnot (!A \land !B)$

Usando as técnicas acima, começamos escrevendo o sequente final desejado
na base.
\begin{prooftree}
\AxiomC{}
\UnaryInf$ \lnot !A \lor \lnot !B \fCenter \lnot (!A \land !B) $
\end{prooftree}
Os conectivos principais disponíveis das !!{sentence}s no sequente final são
o símbolo $\lor$ e o símbolo $\lnot$. Funcionaria aplicar
aqui ou a regra \LeftR{\lor} ou a regra \RightR{\lnot}, mas começamos
com a regra \RightR{\lnot} porque ela evita, por ora, a divisão em dois
ramos:
\begin{prooftree}
\AxiomC{}
\UnaryInf$!A \land !B, \lnot !A \lor \lnot !B \fCenter $
\RightLabel{\RightR{\lnot}}
\UnaryInf$\lnot !A \lor \lnot !B \fCenter \lnot (!A \land !B)$
\end{prooftree}
Agora temos a escolha de olhar para a regra \LeftR{\land} ou a regra
\LeftR{\lor}. Vejamos o que acontece quando aplicamos a regra \LeftR{\land}:
temos a escolha de começar com o sequente $!A,
\lnot !A \lor !B \Sequent \quad$ ou com o sequente $!B, \lnot !A
\lor !B \Sequent \quad$. Como a !!{derivation} é simétrica em
relação a $!A$ e $!B$, vamos com o primeiro:
\begin{prooftree}
\AxiomC{}
\UnaryInf$!A, \lnot !A \lor \lnot !B \fCenter $
\RightLabel{\LeftR{\land}}
\UnaryInf$!A \land !B, \lnot !A \lor \lnot !B \fCenter $
\RightLabel{\RightR{\lnot}}
\UnaryInf$\lnot !A \lor \lnot !B \fCenter \lnot (!A \land !B)$
\end{prooftree}
Continuando a preencher a !!{derivation}, vemos que esbarramos em um problema:
\begin{prooftree}
\Axiom$!A \fCenter !A$
\RightLabel{\LeftR{\lnot}}
\UnaryInf$ \lnot !A, !A \fCenter$
\AxiomC{}
\RightLabel{?}
\UnaryInf$!A \fCenter !B$
\RightLabel{\LeftR{\lnot}}
\UnaryInf$ \lnot !B, !A \fCenter$
\RightLabel{\LeftR{\lor}}
\BinaryInf$\lnot !A \lor \lnot !B, !A \fCenter $
\RightLabel{\LeftR{\Exchange}}
\UnaryInf$!A, \lnot !A \lor \lnot !B \fCenter $
\RightLabel{\LeftR{\land}}
\UnaryInf$!A \land !B, \lnot !A \lor \lnot !B \fCenter $
\RightLabel{\RightR{\lnot}}
\UnaryInf$\lnot !A \lor \lnot !B \fCenter \lnot (!A \land !B)$
\end{prooftree}
O topo do ramo direito não pode ser reduzido mais, e ele
não pode ser levado, por meio de inferências estruturais, a um sequente
inicial, então este não é o caminho certo a tomar. Assim, claramente, foi um
erro aplicar a regra \LeftR{\land} acima. Voltando ao que
tínhamos antes e executando a regra \LeftR{\lor} em vez disso, obtemos
\begin{prooftree}
\AxiomC{}
\UnaryInf$\lnot !A, !A \land !B \fCenter $

\AxiomC{}
\UnaryInf$\lnot !B, !A \land !B \fCenter $

\RightLabel{\LeftR{\lor}}
\BinaryInf$\lnot !A \lor \lnot !B, !A \land !B \fCenter $
\RightLabel{\LeftR{\Exchange}}
\UnaryInf$!A \land !B, \lnot !A \lor \lnot !B \fCenter $
\RightLabel{\RightR{\lnot}}
\UnaryInf$\lnot !A \lor \lnot !B \fCenter \lnot (!A \land !B)$
\end{prooftree}
Completando cada ramo como fizemos antes, obtemos
\begin{prooftree}
\Axiom$ !A \fCenter!A$
\RightLabel{\LeftR{\land}}
\UnaryInf$!A \land !B \fCenter !A$
\RightLabel{\LeftR{\lnot}}
\UnaryInf$\lnot !A, !A \land !B \fCenter $

\Axiom$ !B \fCenter !B$
\RightLabel{\LeftR{\land}}
\UnaryInf$!A \land !B \fCenter !B$
\RightLabel{\LeftR{\lnot}}
\UnaryInf$\lnot !B, !A \land !B \fCenter $

\RightLabel{\LeftR{\lor}}
\BinaryInf$\lnot !A \lor \lnot !B, !A \land !B \fCenter $
\RightLabel{\LeftR{\Exchange}}
\UnaryInf$!A \land !B, \lnot !A \lor \lnot !B \fCenter $
\RightLabel{\RightR{\lnot}}
\UnaryInf$\lnot !A \lor \lnot !B \fCenter \lnot (!A \land !B)$
\end{prooftree}
(Poderíamos ter executado as regras de $\land$ mais abaixo do que as
regras de $\lnot$ nesses passos e ainda assim obtido uma !!{derivation}
correta).
\end{ex}

\begin{ex}
Até agora não usamos a regra de contração, mas às vezes ela é
necessária. Eis um exemplo em que isso acontece.  Suponha que queiramos
provar $\quad \Sequent !A \lor \lnot !A$. Aplicar $\RightR{\lor}$
de trás para frente nos daria uma destas duas !!{derivation}s:
\begin{prooftree}
\AxiomC{}
\UnaryInf$ \fCenter !A$
\RightLabel{\RightR{\lor}}
\UnaryInf$ \fCenter !A \lor \lnot !A$
\DisplayProof\qquad\bottomAlignProof
\AxiomC{}
\UnaryInf$!A \fCenter $
\RightLabel{\RightR{\lnot}}
\UnaryInf$ \fCenter \lnot !A$
\RightLabel{\RightR{\lor}}
\UnaryInf$ \fCenter !A \lor \lnot !A$
\end{prooftree}
Nenhuma destas, é claro, termina em um sequente inicial.  O truque é
perceber que a regra de contração nos permite combinar duas cópias
de !!a{sentence} em uma só---e quando estamos buscando uma prova,
isto é, indo de baixo para cima, podemos manter uma cópia de $!A \lor \lnot
!A$ na premissa, por exemplo,
\begin{prooftree}
\AxiomC{}
\UnaryInf$ \fCenter !A \lor \lnot !A, !A$
\RightLabel{\RightR{\lor}}
\UnaryInf$ \fCenter !A \lor \lnot !A, !A \lor \lnot !A$
\RightLabel{\RightR{\Contraction}}
\UnaryInf$ \fCenter !A \lor \lnot !A$
\end{prooftree}
Agora podemos aplicar $\RightR{\lor}$ uma segunda vez, e também obter~$\lnot
!A$, o que leva a uma !!{derivation} completa.
\begin{prooftree}
\Axiom$!A \fCenter !A$
\RightLabel{\RightR{\lnot}}
\UnaryInf$\fCenter !A, \lnot !A$
\RightLabel{\RightR{\lor}}
\UnaryInf$\fCenter !A, !A \lor \lnot !A$
\RightLabel{\RightR{\Exchange}}
\UnaryInf$ \fCenter !A \lor \lnot !A, !A$
\RightLabel{\RightR{\lor}}
\UnaryInf$ \fCenter !A \lor \lnot !A, !A \lor \lnot !A$
\RightLabel{\RightR{\Contraction}}
\UnaryInf$ \fCenter !A \lor \lnot !A$
\end{prooftree}
\end{ex}

\begin{prob}
Dê !!{derivation}s dos seguintes sequentes:
\begin{enumerate}
\item $!A \land (!B \land !C) \Sequent (!A \land !B) \land !C$.
\item $!A \lor (!B \lor !C) \Sequent (!A \lor !B) \lor !C$.
\item $!A \lif (!B \lif !C) \Sequent !B \lif (!A \lif !C)$.
\item $!A \Sequent \lnot\lnot !A$.
\end{enumerate}
\end{prob}

\begin{prob}
Dê !!{derivation}s dos seguintes sequentes:
\begin{enumerate}
\item $(!A \lor !B) \lif !C \Sequent !A \lif !C$.
\item $(!A \lif !C) \land (!B \lif !C) \Sequent (!A \lor !B) \lif !C$.
\item $\Sequent \lnot(!A \land \lnot !A)$.
\item $!B \lif !A \Sequent \lnot !A \lif \lnot !B$.
\item $\Sequent (!A \lif \lnot !A) \lif \lnot !A$.
\item $\Sequent \lnot(!A \lif !B) \lif \lnot !B$.
\item $!A \lif !C \Sequent \lnot (!A \land \lnot !C)$.
\item $!A \land \lnot !C \Sequent \lnot (!A \lif !C)$.
\item $!A \lor !B, \lnot !B \Sequent !A$.
\item $\lnot !A \lor \lnot !B \Sequent \lnot(!A \land !B)$.
\item $\Sequent (\lnot !A \land \lnot !B) \lif\lnot(!A \lor !B)$.
\item $\Sequent \lnot(!A \lor !B) \lif (\lnot !A \land \lnot !B)$.
\end{enumerate}
\end{prob}

\begin{prob}
Dê !!{derivation}s dos seguintes sequentes:
\begin{enumerate}
\item $\lnot(!A \lif !B) \Sequent !A$.
\item $\lnot(!A \land !B) \Sequent \lnot !A \lor \lnot !B$.
\item $!A \lif !B \Sequent \lnot !A \lor !B$.
\item $\Sequent \lnot \lnot !A \lif !A$.
\item $!A \lif !B, \lnot !A \lif !B \Sequent !B$.
\item $(!A \land !B) \lif !C \Sequent (!A \lif !C) \lor (!B \lif !C)$.
\item $(!A \lif !B) \lif !A \Sequent !A$.
\item $\Sequent (!A \lif !B) \lor (!B \lif !C)$.
\end{enumerate}
(Todos esses requerem a regra $\RightR{\Contraction}$.)
\end{prob}
\end{document}
